% Section 09 - Cas d'usage pratiques
\section{Cas d'usage pratiques et scénarios}

\subsection{Scénario 1 : Fin de semestre - Relevés pour une classe de 150 étudiants}

\textbf{Contexte :}
\begin{itemize}[leftmargin=*]
    \item Classe : Licence 1 Médecine
    \item Semestre : S1
    \item Nombre d'étudiants : 150
    \item Objectif : Générer tous les relevés pour distribution
\end{itemize}

\textbf{Workflow détaillé :}

\begin{enumerate}[leftmargin=*]
    \item \textbf{Préparation} (15 min)
    \begin{itemize}
        \item Créer la configuration "L1 Médecine 2023-2024" dans \menu{Config des relevés}
        \item Ajouter le semestre S1 avec toutes les UEs et ECs
        \item Configurer les coefficients et crédits
        \item Personnaliser le thème (optionnel)
    \end{itemize}

    \item \textbf{Préparation du fichier Excel} (10 min)
    \begin{itemize}
        \item Vérifier toutes les colonnes obligatoires présentes
        \item Vérifier que les noms d'ECs correspondent exactement à la configuration
        \item Enregistrer au format .xlsx
    \end{itemize}

    \item \textbf{Génération} (5-10 min)
    \begin{itemize}
        \item Charger le fichier Excel (150 lignes)
        \item Sélectionner la configuration "L1 Médecine 2023-2024"
        \item Choisir le semestre S1
        \item Mapper les colonnes (mémorisé pour la prochaine fois)
        \item Format d'export : ZIP avec compression
        \item QR code chiffré : Activé
        \item Lancer la génération
    \end{itemize}

    \item \textbf{Vérification} (5 min)
    \begin{itemize}
        \item Ouvrir 5-10 relevés au hasard
        \item Vérifier les notes, moyennes, mentions
        \item Tester un QR code
    \end{itemize}

    \item \textbf{Distribution}
    \begin{itemize}
        \item Extraire le ZIP (150 fichiers PDF)
        \item Imprimer ou distribuer numériquement
    \end{itemize}
\end{enumerate}

\textbf{Résultat attendu :}
\begin{itemize}[leftmargin=*]
    \item 150 relevés PDF individuels dans une archive ZIP
    \item Taille totale : environ 30-40 Mo (avec compression)
    \item Temps total : 35-40 minutes
    \item Chaque relevé contient toutes les notes, moyennes calculées, mention, et QR code
\end{itemize}

\subsection{Scénario 2 : Génération multi-semestres pour relevé de parcours}

\textbf{Contexte :}
\begin{itemize}[leftmargin=*]
    \item Étudiant finissant sa Licence (L3)
    \item Besoin : Relevé complet de L1 à L3 (6 semestres)
    \item Pour : Dossier de candidature en Master
\end{itemize}

\textbf{Workflow :}

\begin{enumerate}[leftmargin=*]
    \item Charger le fichier Excel contenant TOUS les semestres
    \item Sélectionner la configuration de classe
    \item Choisir \textbf{"Tous les semestres"}
    \item Cocher S1, S2, S3, S4, S5, S6
    \item Format d'export : \textbf{PDF unique} (tout fusionné en un fichier)
    \item Générer
\end{enumerate}

\textbf{Résultat :}
\begin{itemize}[leftmargin=*]
    \item 1 fichier PDF de 6 pages (1 page par semestre)
    \item Nom : \texttt{Releves\_Fusionnes\_DUPONT\_Jean\_2024.pdf}
    \item Parfait pour candidatures et archivage
\end{itemize}

\subsection{Scénario 3 : Attestations de réussite en fin d'année}

\textbf{Contexte :}
\begin{itemize}[leftmargin=*]
    \item Fin de l'année académique
    \item 200 étudiants de L1
    \item Besoin : Attestations pour les étudiants admis uniquement
\end{itemize}

\textbf{Workflow :}

\begin{enumerate}[leftmargin=*]
    \item Menu \menu{Générer les attestations}
    \item Charger le fichier Excel avec toutes les notes et \texttt{MOYENNE\_GENERALE}
    \item L'application affiche automatiquement :
    \begin{itemize}
        \item 167 étudiants ADMIS (moyenne $\geq$ 10)
        \item 33 étudiants AJOURNÉS (exclus automatiquement)
        \item Taux de réussite : 83.5\%
    \end{itemize}
    \item Choisir un préréglage de thème : "Professionnel Académique"
    \item Activer le QR code chiffré
    \item Format : ZIP avec compression
    \item Générer
\end{enumerate}

\textbf{Résultat :}
\begin{itemize}[leftmargin=*]
    \item 167 attestations PDF (uniquement les admis)
    \item Archive ZIP de 15-20 Mo
    \item Temps de génération : 8-10 minutes
\end{itemize}

\subsection{Scénario 4 : Ajout de QR codes sur relevés déjà imprimés}

\textbf{Contexte :}
\begin{itemize}[leftmargin=*]
    \item 100 relevés déjà générés et distribués SANS QR code
    \item Nouvelle politique : Tous les relevés doivent avoir un QR de vérification
    \item Documents déjà scannés en PDF
\end{itemize}

\textbf{Workflow :}

\begin{enumerate}[leftmargin=*]
    \item Menu \menu{QR Codes sur PDF} (nécessite licence active)
    \item Charger les 100 PDFs scannés
    \item Charger le fichier Excel avec les données étudiants
    \item Vérifier le matching automatique (devrait être 100\% si noms de fichiers contiennent matricules)
    \item Corriger manuellement les non-matchés
    \item Positionner le QR : Haut Droite (85\%, 15\%)
    \item Colonnes dans QR : MATRICULE, NOM, PRENOM, SEMESTRE
    \item Taille : Moyen, Correction d'erreur : M
    \item Format : ZIP
    \item Traiter tous les documents
\end{enumerate}

\textbf{Résultat :}
\begin{itemize}[leftmargin=*]
    \item 100 PDFs avec QR codes ajoutés
    \item QR présent sur toutes les pages de chaque document
    \item Prêts à être réimprimés
\end{itemize}

\subsection{Scénario 5 : Création d'attestations bilingues (FR/EN)}

\textbf{Contexte :}
\begin{itemize}[leftmargin=*]
    \item Étudiants internationaux demandant des attestations en anglais
    \item 50 étudiants admis
    \item Besoin : Attestations en anglais
\end{itemize}

\textbf{Préparation du fichier Excel :}

Ajouter les colonnes de traduction :
\begin{itemize}[leftmargin=*]
    \item \texttt{DOMAINE\_EN} : "Health Sciences"
    \item \texttt{PARCOURS\_EN} : "General Medicine"
    \item \texttt{SPECIALITE\_EN} : "Cardiology"
    \item \texttt{MENTION\_EN} : "Excellent", "Good", "Fair", "Pass"
\end{itemize}

\textbf{Workflow :}

\begin{enumerate}[leftmargin=*]
    \item Charger l'Excel avec colonnes FR et EN
    \item Vérifier l'éligibilité (50 admis détectés)
    \item Dans l'onglet \textbf{Paramètres}, sélectionner \champ{Langue : Anglais}
    \item Choisir un thème approprié
    \item Générer
\end{enumerate}

\textbf{Résultat :}
\begin{itemize}[leftmargin=*]
    \item 50 attestations en anglais
    \item Titre : "Certificate of Achievement"
    \item Toutes les valeurs traduites
\end{itemize}

\subsection{Scénario 6 : Test et configuration initiale}

\textbf{Contexte :}
\begin{itemize}[leftmargin=*]
    \item Première utilisation de l'application
    \item Besoin : Configurer et tester avant utilisation en production
\end{itemize}

\textbf{Workflow recommandé :}

\begin{enumerate}[leftmargin=*]
    \item \textbf{Jour 1 : Configuration de base}
    \begin{itemize}
        \item Menu \menu{Configurer les entêtes}
        \item Renseigner toutes les informations de l'établissement
        \item Charger le logo
        \item Enregistrer
    \end{itemize}

    \item \textbf{Jour 2 : Créer une configuration de test}
    \begin{itemize}
        \item Menu \menu{Config des relevés}
        \item Créer une configuration simplifiée : "Test - L1"
        \item 1 semestre, 2 UEs, 3-4 ECs par UE
        \item Enregistrer
    \end{itemize}

    \item \textbf{Jour 3 : Test avec données fictives}
    \begin{itemize}
        \item Créer un fichier Excel avec 5-10 lignes de données fictives
        \item Générer des relevés de test
        \item Vérifier le résultat
        \item Ajuster le thème si nécessaire
    \end{itemize}

    \item \textbf{Jour 4 : Test d'attestations}
    \begin{itemize}
        \item Ajouter colonne \texttt{MOYENNE\_GENERALE} à l'Excel de test
        \item Générer des attestations de test
        \item Tester les différents thèmes prédéfinis
        \item Choisir celui qui convient
    \end{itemize}

    \item \textbf{Jour 5 : Configuration réelle}
    \begin{itemize}
        \item Créer les vraies configurations (toutes les classes)
        \item Personnaliser les thèmes définitifs
        \item Exporter les configurations pour sauvegarde
        \item Prêt pour production
    \end{itemize}
\end{enumerate}

\subsection{Scénario 7 : Gestion d'une promotion complète (L1 à L3)}

\textbf{Contexte :}
\begin{itemize}[leftmargin=*]
    \item Faculté avec 3 niveaux (L1, L2, L3)
    \item Environ 500 étudiants au total
    \item Fin de semestre 2
\end{itemize}

\textbf{Stratégie organisée :}

\begin{enumerate}[leftmargin=*]
    \item \textbf{Créer 3 configurations séparées}
    \begin{itemize}
        \item Config L1 (S1 et S2)
        \item Config L2 (S3 et S4)
        \item Config L3 (S5 et S6)
    \end{itemize}

    \item \textbf{Préparer 3 fichiers Excel}
    \begin{itemize}
        \item Excel\_L1.xlsx (150 étudiants)
        \item Excel\_L2.xlsx (180 étudiants)
        \item Excel\_L3.xlsx (170 étudiants)
    \end{itemize}

    \item \textbf{Générer par niveau}
    \begin{itemize}
        \item Jour 1 : Relevés L1 pour S2
        \item Jour 2 : Relevés L2 pour S4
        \item Jour 3 : Relevés L3 pour S6
        \item Vérification quotidienne des résultats
    \end{itemize}

    \item \textbf{Attestations pour les finissants L3}
    \begin{itemize}
        \item Générer attestations uniquement pour L3 admis
        \item Thème plus solennel
        \item QR code obligatoire
    \end{itemize}
\end{enumerate}

\textbf{Avantages de cette approche :}
\begin{itemize}[leftmargin=*]
    \item Évite la surcharge mémoire
    \item Permet des vérifications intermédiaires
    \item Facilite la correction d'erreurs
    \item Organisation claire par niveau
\end{itemize}

\subsection{Conseils généraux pour l'utilisation quotidienne}

\begin{enumerate}[leftmargin=*]
    \item \textbf{Gardez vos fichiers Excel organisés}
    \begin{itemize}
        \item Nommage cohérent : \texttt{Notes\_L1\_S1\_2023-2024.xlsx}
        \item Dossier dédié par année académique
        \item Sauvegardes régulières
    \end{itemize}

    \item \textbf{Exportez vos configurations régulièrement}
    \begin{itemize}
        \item Sauvegarde mensuelle des configurations
        \item Stockage sur cloud ou disque externe
        \item Facilite le partage entre collègues
    \end{itemize}

    \item \textbf{Testez toujours avec un petit échantillon}
    \begin{itemize}
        \item Avant une génération massive, testez avec 5-10 étudiants
        \item Vérifiez le résultat
        \item Puis lancez la génération complète
    \end{itemize}

    \item \textbf{Utilisez l'historique}
    \begin{itemize}
        \item Consultez régulièrement \menu{Historique des docs}
        \item Tracez toutes vos générations
        \item Exportez l'historique trimestriellement
    \end{itemize}

    \item \textbf{Standardisez vos thèmes}
    \begin{itemize}
        \item Créez un thème officiel pour chaque type de document
        \item Exportez-les
        \item Partagez avec vos collègues pour cohérence
    \end{itemize}
\end{enumerate}

\newpage
